% !TEX root = ../main.tex
% =====================================================================
% 7. CONCLUSION AND FUTURE WORK
%
% No new information. The full objective table is in Appendix A and is
% referred to here rather than repeated.
% =====================================================================

\section{Conclusion}
\label{sec:conclusion}

This internship addressed a question that arises whenever a time series classifier has to run on a
constrained device: can such a model be made smaller \emph{and} more interpretable at the same time,
rather than trading one against the other? The starting point was \millet{} \citep{early2024millet},
which produces a faithful built-in explanation but needs about 424\,K parameters.

The proposed solution is \seanet{}, an architecture in two interchangeable parts: an encoder built
from multi-scale depthwise-separable dilated convolutional blocks that never collapse the time axis,
in five variants plus two wrappers (\Cref{sec:methods:encoder}), and seven MIL pooling heads
(\Cref{sec:methods:pooling}), each addressing a specific limitation of \millet{}'s conjunctive
pooling and five of them provably reducing to it. Seventy-two combinations were evaluated on
\webtraffic{} and the \ucr{} archive against FCN, ResNet and \millet{} baselines re-trained under an
identical training configuration.

The main finding is that the three requirements of \Cref{sec:intro:problem} are compatible. On
\webtraffic{}, \mbot{} exceeds the re-trained \millet{} baseline on accuracy, AOPCR and NDCG@$n$
simultaneously while using 90\,\% fewer parameters (41{,}324 against 423{,}707), a third of the file
size and less time per prediction; \mgat{} is the most accurate model of the study at 0.955 with
61{,}740 parameters and the smallest seed-to-seed spread measured. Two results explain \emph{why}:
the parameter reduction targets channel mixing rather than temporal filtering, and the explanation
gain comes from Top-$k$ selection making the interpretation exactly zero where the model used no
evidence -- both supported by controlled ablations (\Cref{sec:discussion}). A third finding is that
the encoder and the pooling head interact, so the best model is not obtained by combining the two
individually best components.

The principal limitation is equally clear and constrains the claim above. The headline models were
selected on \webtraffic{}, the only dataset with per-time-step ground truth, and over the 85 \ucr{}
datasets the re-trained \millet{} baseline remains ahead on average (\Cref{sec:res:ucr}). A choice
made on one dataset should not be expected to be optimal on 85 heterogeneous ones, and it was not.
The remaining limitations are listed in \Cref{sec:disc:limitations}. The objectives set at the start
of the internship and their completion status are summarised in \Cref{tab:objectives}: five of the
six were reached, and the sixth -- deployment on a physical microcontroller -- was not attempted.

% ---------------------------------------------------------------------
\subsection{Future work}
\label{sec:future}

\begin{itemize}\itemsep3pt
  \item \textbf{Re-train \seanet{} under the long training configuration.} \millet{}'s own
        configuration recovers the published \ucr{} accuracy on this pipeline
        (\Cref{tab:ucr_accuracy}), so applying the same four-times-longer budget to the \seanet{}
        encoders would determine whether the \ucr{} gap of \Cref{sec:res:ucr} is architectural or
        merely a budget effect. This is the single most informative experiment left.

  \item \textbf{Adapt further MIL pooling methods from the histopathology literature.} The two heads
        transferred here (\Cref{sec:pool:gated}, \Cref{sec:pool:dsmil}) come from that literature,
        and two more recent methods address exactly the failure mode observed in
        \Cref{sec:disc:negative}. DTFD-MIL \citep{zhang2022dtfdmil} splits a bag into pseudo-bags,
        trains a first-tier model on them and distils the resulting instance features into a
        second tier; this increases the number of effective training bags, which is a plausible
        remedy for a head whose attention is unstable early in training. ReMix
        \citep{yang2022remix} reduces a bag to a small set of representative instance prototypes and
        augments it by mixing them, which addresses the dilution of \Cref{sec:bg:millet} at the level
        of the bag rather than the aggregator. Neither has been implemented here; both are proposed
        as the next pooling heads to try, with the caveat that whole-slide images contain thousands
        of instances per bag while these series contain around a thousand time steps, so the
        transfer is not automatic.

  \item \textbf{Make the dilation cap depend on the series length.} Over all 128 \ucr{} datasets,
        accuracy falls with series length (correlation $-0.34$; 0.857 on the short half of the
        archive against 0.755 on the long half), which is what a fixed cap of 16 predicts and an
        effect far larger than any seed-to-seed spread measured (\Cref{tab:seeds}). Adapting the cap
        is a small change with a clear hypothesis behind it.

  \item \textbf{Correct the two heads that failed for optimisation reasons.} Initialising the
        dual-stream blend $\lambda$ near zero, or penalising it away from the safe branch early in
        training, is a direct test of the diagnosis in \Cref{sec:disc:negative}.

  \item \textbf{Evaluate on a real dataset with per-time-step annotations,} so that explanation
        correctness no longer rests on a generator's own record of what it modified.

  \item \textbf{Deploy on the device.} Quantise \mbot{} \citep{jacob2018quantization} and measure
        accuracy, latency and memory on a physical microcontroller, which is what objective~O6
        requires and the natural continuation of this project.
\end{itemize}

\subsubsection*{Acknowledgements}

I thank my supervisor, Tanmoy Mondal, for his guidance throughout this internship, and the Fox team
for their support. I also thank the Grid'5000 testbed \citep{balouek2013grid5000}, whose GPU nodes
made the 72-model study in this report possible.
